\CWHeader{Лабораторная работа \textnumero 5}

\CWProblem{
Необходимо реализовать поиск образцов в заранее известном тексте с использованием суффиксного массива.

На вход программе подаётся текст, расположенный на первой строке. Далее, до конца файла, следуют строки с образцами, которые необходимо найти в этом тексте.

Для каждого найденного образца требуется вывести строку, начинающуюся с последовательного номера данного образца и двоеточия. После двоеточия через запятую должны быть перечислены номера позиций, в которых данный образец встречается в тексте. Позиции необходимо выводить в порядке возрастания.

Ограничения задачи:
\begin{itemize}
    \item ограничение времени: 1 секунда;
    \item ограничение памяти: 1280 Мб;
    \item ввод: стандартный ввод или \texttt{input.txt};
    \item вывод: стандартный вывод или \texttt{output.txt}.
\end{itemize}

Пример входных данных:
\begin{alltt}
abcdabc
abcd
bcd
bc
\end{alltt}

Пример выходных данных:
\begin{alltt}
1: 1
2: 2
3: 2, 6
\end{alltt}

В работе используется подход на основе суффиксного массива. Для ускорения поиска дополнительно применяются массив LCP и структура RMQ, позволяющая быстро находить длину наибольшего общего префикса двух суффиксов. Это позволяет уменьшить количество повторных сравнений символов при бинарном поиске образца в суффиксном массиве.

Результатом лабораторной работы является программа, которая:
\begin{itemize}
    \item строит суффиксный массив для заданного текста;
    \item строит массив LCP;
    \item строит структуру RMQ для быстрых запросов минимума на отрезке массива LCP;
    \item для каждого образца находит диапазон суффиксов, начинающихся с данного образца;
    \item выводит все позиции вхождений образца в порядке возрастания.
\end{itemize}
}

\pagebreak